\chapter{TESTING AND VERIFICATION}

\section{Test Cases}
Testing was performed through both manual execution and automated regression
tests. Manual runs were used to inspect the full pipeline from source code to
output, while automated tests were added for the semantic and runtime behaviour
of the implemented type checker and interpreter.

\begin{table}[H]
\centering
\caption{Representative Test Cases}
\begin{tabular}{p{3.5cm}p{5cm}p{4cm}}
\toprule
Category & Program Behaviour & Expected Result \\
\midrule
Arithmetic & Evaluate mixed numeric expressions & Correct value and inferred type \\
If-then-else & Choose one branch from a Boolean condition & Only matching branch runs \\
While-loop & Print during repeated execution & One output line per iteration \\
Functions & Infer parameters and return type & Valid call and correct return value \\
Type mismatch & Reassign incompatible type & Type error before runtime \\
\bottomrule
\end{tabular}
\end{table}

\subsection{Lexer and Symbol Table}
The lexer was verified manually by running the full pipeline and inspecting
Stage~1 token output for programs that exercise every token category: integer
and float literals, Boolean literals, string literals, all six keywords,
two-character comparison operators (\texttt{==}, \texttt{!=}), single-character
operators, and delimiters. Identifier scanning with keyword fallback was
confirmed by checking that \texttt{true}, \texttt{false}, and reserved words
are emitted with their keyword token types while non-keyword names produce
\texttt{IDENTIFIER} tokens.

Error paths were also checked: supplying a bare \texttt{!} character raises a
\texttt{LexerError} with a source position and a message suggesting
\texttt{!=}, and an unterminated string literal raises a \texttt{LexerError}
citing the opening line.

The symbol table was verified through the type checker and the Stage~3 output.
Duplicate variable definitions were confirmed to raise a
\texttt{SymbolTableError} before execution. Scope isolation was verified with
nested blocks: a variable defined inside an \texttt{if} or \texttt{while} body
is not visible in the outer scope after the block ends. The
\texttt{format\_table()} output was checked against the expected column layout
(name, kind, type, initialisation status, parameters) for both variable and
function symbols.

\subsection{Parser and AST}
The parser was verified manually by running the full pipeline and inspecting
Stage~2 AST output. The following cases were checked:

\begin{itemize}
    \item arithmetic expressions with mixed precedence (e.g.\ \texttt{x + y * 2}) produce a tree where multiplication is a deeper sub-tree than addition, confirming the \texttt{\_expr}/\texttt{\_term}/\texttt{\_factor} precedence encoding,
    \item assignment versus expression statement dispatch: \texttt{x = 5;} produces an \texttt{Assign} node while \texttt{add(1,2);} produces a \texttt{FunctionCall} node at the statement level, confirming the IDENTIFIER+ASSIGN lookahead,
    \item \texttt{if} with and without an \texttt{else} clause: the resulting \texttt{If} node has a non-null \texttt{else\_block} only when the clause is present,
    \item function definitions store only parameter names in the \texttt{FunctionDef} node; the AST printer output was inspected to confirm no type annotations are present at this stage,
    \item string literals in \texttt{Literal} nodes have their surrounding double-quotes stripped; \texttt{print("hello")} prints \texttt{hello : String} at runtime,
    \item parse errors: missing semicolons, unmatched parentheses, and unexpected tokens all produce \texttt{[line X, col Y] ParseError:} messages with correct source positions.
\end{itemize}

\subsection{Arithmetic Expressions}
Arithmetic expressions were tested using both integer-only and mixed
integer-float programs. These tests verified precedence and numeric promotion.
Expressions such as \texttt{x + y * 2} confirmed that multiplication is applied
before addition.

\subsection{Boolean Expressions}
Boolean expressions were tested inside \texttt{if} and \texttt{while}
conditions. The checker verified valid comparison operands, and the interpreter
used the resulting Boolean values for control flow.

\subsection{Assignment Statements}
Assignments were tested for both first-use inference and reassignment checks.
For example, a variable assigned \texttt{5} and later assigned
\texttt{"hello"} correctly triggers a type error.

\subsection{If-Then-Else}
Conditional programs were used to verify that exactly one branch is executed and
that non-Boolean conditions are rejected before runtime.

\subsection{While-Loop}
While-loops were tested with a countdown example. This confirmed that the loop
body executes repeatedly and that \texttt{print()} emits output progressively on
each iteration rather than only once at the end.

\subsection{Functions}
Function tests covered parameter passing, local execution, and return values.
The implementation infers parameter types from the first call and then enforces
that signature for later calls.

\subsection{Print}
The built-in \texttt{print()} function was tested with all supported runtime
types. The output includes the printed value together with its type, such as
\texttt{29.5 : Float} and \texttt{"plc" : String}.

\subsection{Unary Minus}
The unary minus operator was verified for integer literals (\texttt{x = -5}),
float literals (\texttt{y = -3.14}), negated variable references
(\texttt{y = -x}), and negation inside a larger expression
(\texttt{x = 3 + -2}). In each case the type checker assigned the correct
type and the interpreter produced the expected value.

\section{Type Error Detection}
The checker was verified with invalid programs involving incompatible
assignment, invalid arithmetic operands, wrong function arguments, and
non-Boolean conditions. In each case the system stopped before execution and
reported a source-positioned type error. Representative error messages are
shown below.

\begin{lstlisting}
# Assignment type mismatch
[line 2, col 1] TypeError: Cannot assign String to variable 'x' of type Integer

# Non-Boolean if condition
[line 1, col 4] TypeError: If condition must have type Boolean

# Wrong argument type for function
[line 5, col 9] TypeError: Argument 1 of 'add' must be Integer, got String
\end{lstlisting}

\section{Automated Test Suite}
Automated regression tests were added in
\texttt{tests/test\_member3\_pipeline.py} using the \texttt{unittest}
framework. The suite currently covers:

\begin{itemize}
	\item progressive output during while-loop execution,
	\item assignment type mismatch detection,
	\item function parameter and return type inference.
\end{itemize}

The tests can be run using:

\begin{lstlisting}
python3 -m unittest discover -s tests -v
\end{lstlisting}

All seven tests pass. The four unary-minus tests were added when support for
negation was introduced, confirming that the feature is covered alongside the
original three semantic checks.

\section{Error Handling}
The language system reports errors at three main stages:

\begin{itemize}
	\item lexical errors for malformed characters or unterminated strings,
	\item parse errors for invalid syntax such as missing semicolons,
	\item type errors for programs that violate static typing rules.
\end{itemize}

Runtime errors are also reported with source positions whenever execution
reaches an invalid state such as an undefined variable reference.
